\chapter*{ВВЕДЕНИЕ}
\addcontentsline{toc}{chapter}{ВВЕДЕНИЕ}
\sloppy

Машинное обучение за последнее десятилетие стало одной из ключевых компетенций для IT-специалистов. Спрос на образовательные программы по данному направлению неуклонно растет: курсы предлагают ведущие университеты (НИУ ВШЭ, МГУ, МФТИ), корпоративные школы (ШАД Яндекса, Т-Образование) и открытые платформы (ODS, Stepik, Karpov.Courses). Вместе с ростом числа слушателей обостряется проблема масштабируемости оценивания: преподаватели тратят значительную часть времени на проверку работ, а при увеличении потока студентов качество обратной связи снижается.

Анализ существующих курсов по машинному обучению показал, что автоматизация оценивания остается фрагментарной. Существующие инструменты - юнит-тесты в Jupyter Notebooks, контестные системы, тесты с выбором ответа - эффективны для проверки кода и формализованных ответов, но не позволяют оценивать задания открытого типа: теоретические вопросы, развернутые объяснения, интерпретацию результатов. Такие задания по-прежнему проверяются вручную, что создает узкое место в образовательном процессе. Появление больших языковых моделей (БЯМ) открывает возможность автоматизировать проверку заданий, требующих понимания естественного языка и оценки содержательной корректности ответов.

Новизна работы заключается в разработке и апробации формата образовательного курса по машинному обучению, в котором автоматическая оценка открытых ответов реализована с использованием БЯМ, интегрированного в существующий учебный пайплайн nbgrader. В отличие от решений, ориентированных только на проверку кода и закрытых тестов, предложенный подход объединяет новый формат курса, встроенный БЯМ и его применение на реальной аудитории студентов.

Практическая значимость работы определяется внедрением разработанного курса в образовательный процесс НИУ ВШЭ, а также апробацией системы автоматической проверки на реальной аудитории студентов, что позволяет оценить её применимость и эффективность в условиях учебного процесса и потенциальную переносимость на другие технические дисциплины.

Целью данной работы является разработка автоматизированной системы проверки учебных заданий на основе больших языковых моделей для курса по машинному обучению, ее интеграция в образовательный процесс и апробация.

Для достижения данной цели были поставлены следующие задачи:
\begin{enumerate}
  \item Провести анализ существующих курсов по машинному обучению и сформировать профиль целевой аудитории.
  \item Исследовать подходы к автоматизированной проверке учебных заданий в образовательной среде.
  \item Разработать методологию преподавания и структуру учебного курса по машинному обучению.
  \item Разработать контрольные задания, критерии оценивания и учебно-методические материалы курса.
  \item Реализовать автоматизированную систему проверки учебных заданий с использованием больших языковых моделей.
  \item Провести апробацию разработанного курса и системы автоматической проверки в образовательном процессе НИУ ВШЭ и проанализировать полученные результаты.
\end{enumerate}

Критериями успешности решения выступают: сокращение времени проверки учебных заданий, высокая степень согласованности автоматических оценок с экспертными, а также положительная оценка скорости и качества обратной связи со стороны студентов по итогам апробации.

Работа состоит из введения, трех глав и заключения. В первой главе описана предметная область, проведен анализ существующих курсов и подходов к автоматизированной проверке, сформулированы требования к решению. Во второй главе представлена архитектура разработанного курса и системы автоматической проверки, описаны методология, структура заданий и механизм оценивания на основе БЯМ. В третьей главе приведены результаты апробации курса, анализ метрик качества автоматической проверки и обратная связь от студентов.
